\documentclass[10pt,a4paper]{amsart}
\usepackage[margin=19mm]{geometry}
\usepackage{amsmath,amssymb,mathtools}
\usepackage[scr=rsfs,cal=euler]{mathalfa}
\usepackage{xfrac}
\usepackage[unicode,hidelinks,bookmarks=false]{hyperref}
\newcommand{\ie}{i.e.\ }
\newcommand{\cf}{cf.\ }
\newcommand{\resp}{respectively\ }
\newcommand{\Subsection}{\subsection}
\providecommand{\nobreakdash}{\nobreak\hbox{-}}
\setlength{\emergencystretch}{2em}
\multlinegap0pt
\usepackage{amssymb}
\newtheorem{theorem}{Theorem}
\title{On directional second-order tangent sets of analytic sets and applications in optimization}
\author{Le Cong Trinh}
\date{Source: arXiv:2601.09991v2 (CC BY 4.0)}
\begin{document}
\maketitle

\noindent\emph{Source and licence.} On directional second-order tangent sets of analytic sets and applications in optimization. Version arXiv:2601.09991v2 (CC BY 4.0), distributed by arXiv under the Creative Commons Attribution 4.0 International licence, http://creativecommons.org/licenses/by/4.0/, as recorded in the arXiv licence metadata for this exact version and shown on the abstract page https://arxiv.org/abs/2601.09991v2. Full licence notice: \texttt{licenses/arxiv-2601.09991-CC-BY-4.0.txt}.\par\smallskip

\noindent\textit{Editorial scope.} Counted unit: the article's theorem thm:SOSC (source0.tex lines 1816--1860). The article's complete original proof of it is delivered with it (source0.tex lines 1862--1987, one \texttt{proof} environment of 126 lines). No mathematics is written, repaired or re-derived: the statement and the proof are the article's own lines, in the article's own notation, and no retained line is substituted or re-flowed.  No further source-local context is needed: every cross-reference of every retained line resolves inside the two delivered spans.  Nothing was omitted from any retained span, so the package carries no omission record.  No retained line contains a citation, so the wrapper carries no bibliography and no bibliography entry.  No retained line contains a figure environment, an image-inclusion command or drawing code, so no image is delivered and the package declares no asset dependency.  Editorial formatting only, in addition: an AMS \texttt{amsart} preamble that restates the article's own declarations and macros used by the retained lines, namely the environments \texttt{theorem} on separate counters, together with the standard packages those lines need.  The retained environments print in this document as Theorem 1 in order of appearance, whereas the article prints the same items with the numbering of its own full document.  The complete native source of the article as distributed in the arXiv e-print for this exact version is bundled unchanged as \texttt{sources/192730/source0.tex}.
\begin{theorem}[Second-order sufficiency via $T^2_{0,u}X$]
\label{thm:SOSC}
Let $X\subseteq\mathbb K^n$ be a closed set with $0\in X$, and let
$f\in C^2(\mathbb K^n)$. Assume:
\begin{itemize}
\item[(i)] 
\[
\langle \nabla f(0),u\rangle \ge 0
\qquad \forall\,u\in T_0X;
\]

\item[(ii)] for every $u\in T_0X\setminus\{0\}$ satisfying
\[
\langle \nabla f(0),u\rangle =0,
\]
one has
\begin{equation}\label{eq:SOSC-assump}
\inf_{w\in T^2_{0,u}X}
\Bigl(
\langle u,\nabla^2 f(0)u\rangle
+
\langle \nabla f(0),w\rangle
\Bigr)>0;
\end{equation}

\item[(iii)] \emph{(parabolic regularity at $0$)} for every sequence
$\{x_k\}\subset X\setminus\{0\}$ with $x_k\to 0$ and
\[
\frac{f(x_k)-f(0)}{\|x_k\|^2}
\]
bounded above, there exist, after passing to a subsequence,
a nonzero vector $u\in T_0X$, positive numbers $t_k\downarrow 0$, and a bounded
sequence $\{w_k\}\subseteq\mathbb K^n$ such that
\begin{equation}\label{eq:parabolic-expansion}
x_k=t_k u+\frac12 t_k^2w_k+o(t_k^2),
\end{equation}
and every cluster point of $\{w_k\}$ belongs to $T^2_{0,u}X$.
\end{itemize}
Then $0$ is a strict local minimizer of $f$ on $X$ and satisfies quadratic
growth: there exist $c>0$ and $\varepsilon>0$ such that
\[
f(x)\ge f(0)+c\|x\|^2
\qquad \forall\,x\in X\cap B_\varepsilon(0).
\]
\end{theorem}

\begin{proof}
Suppose, to the contrary, that quadratic growth fails. Then there exists a
sequence $\{x_k\}\subseteq X\setminus\{0\}$ such that $x_k\to 0$ and
\begin{equation}\label{eq:no-quadratic-growth}
\frac{f(x_k)-f(0)}{\|x_k\|^2}\le \frac1k
\qquad (\forall k).
\end{equation}
In particular, the quotients
\(
\frac{f(x_k)-f(0)}{\|x_k\|^2}
\)
are bounded above. By assumption {\rm(iii)}, after passing to a subsequence we
may assume that there exist $u\in T_0X\setminus\{0\}$, positive numbers
$t_k\downarrow 0$, and a bounded sequence $\{w_k\}$ such that
\[
x_k=t_k u+\frac12 t_k^2w_k+o(t_k^2),
\]
and every cluster point of $\{w_k\}$ belongs to $T^2_{0,u}X$.

Since $u\neq 0$, it follows from the expansion above that
\(
\frac{\|x_k\|}{t_k}\to \|u\|,
\)
hence there exist constants $c_1,c_2>0$ such that
\begin{equation}\label{eq:norm-comparison}
c_1 t_k \le \|x_k\| \le c_2 t_k
\qquad \text{for all } k \text{ sufficiently large}.
\end{equation}

Using the Taylor expansion of $f$ at $0$, we obtain
\[
f(x_k)=f(0)
+\langle \nabla f(0),x_k\rangle
+\frac12\langle x_k,\nabla^2 f(0)x_k\rangle
+o(\|x_k\|^2).
\]
Substituting
\[
x_k=t_k u+\frac12 t_k^2w_k+o(t_k^2)
\]
and using \eqref{eq:norm-comparison}, we get
\[
f(x_k)=f(0)
+t_k\langle \nabla f(0),u\rangle
+\frac12 t_k^2
\Bigl(
\langle u,\nabla^2 f(0)u\rangle
+
\langle \nabla f(0),w_k\rangle
\Bigr)
+o(t_k^2).
\]

We first claim that
\(
\langle \nabla f(0),u\rangle =0.
\)
Indeed, by assumption {\rm(i)},
\(
\langle \nabla f(0),u\rangle \ge 0.
\)
If this quantity were strictly positive, say
\(
\langle \nabla f(0),u\rangle =\eta>0,
\)
then for $k$ large enough we would have
\(
f(x_k)-f(0)\ge \frac{\eta}{2}t_k,
\)
and therefore, by \eqref{eq:norm-comparison},
\[
\frac{f(x_k)-f(0)}{\|x_k\|^2}
\ge \frac{\eta}{2c_2^2}\frac{1}{t_k}\to +\infty,
\]
contradicting \eqref{eq:no-quadratic-growth}. Hence
\(
\langle \nabla f(0),u\rangle =0.
\)

Now let $w$ be a cluster point of $\{w_k\}$. By assumption {\rm(iii)},
\(
w\in T^2_{0,u}X.
\)
By assumption {\rm(ii)},
\[
\langle u,\nabla^2 f(0)u\rangle + \langle \nabla f(0),w\rangle >0
\qquad \forall\,w\in T^2_{0,u}X.
\]
Since $T^2_{0,u}X$ is closed and the function
\[
w\longmapsto \langle u,\nabla^2 f(0)u\rangle + \langle \nabla f(0),w\rangle
\]
is continuous, the strict positivity of the infimum in \eqref{eq:SOSC-assump}
implies that there exists $\alpha>0$ such that
\[
\langle u,\nabla^2 f(0)u\rangle + \langle \nabla f(0),w\rangle \ge 2\alpha
\quad \forall\,w\in T^2_{0,u}X.
\]
Passing to a subsequence if necessary, we may assume that $w_k\to w$ for some
cluster point $w\in T^2_{0,u}X$. Then
\[
\langle u,\nabla^2 f(0)u\rangle + \langle \nabla f(0),w_k\rangle
\to
\langle u,\nabla^2 f(0)u\rangle + \langle \nabla f(0),w\rangle
\ge 2\alpha.
\]
Hence, for all $k$ sufficiently large,
\[
\langle u,\nabla^2 f(0)u\rangle + \langle \nabla f(0),w_k\rangle \ge \alpha.
\]
Returning to the Taylor expansion, we obtain
\[
f(x_k)-f(0)\ge \frac12\alpha t_k^2+o(t_k^2).
\]
Therefore, for $k$ large,
\[
f(x_k)-f(0)\ge \frac{\alpha}{4}t_k^2.
\]
Using again \eqref{eq:norm-comparison}, we conclude that
\[
f(x_k)-f(0)\ge \frac{\alpha}{4c_2^2}\|x_k\|^2
\]
for all sufficiently large $k$, contradicting \eqref{eq:no-quadratic-growth}.
This contradiction proves quadratic growth. In particular, $0$ is a strict local
minimizer of $f$ on $X$.
\end{proof}

\end{document}
